\subsection{完成化截面的平衡原子失明、零温半熵律与地基 Parry 流的不可逆性}\label{subsec:conclusion-completed-slice-atomblind-groundstate-escort-current}
在小节 \ref{subsec:conclusion-realinput40-primitive-surgery-residue-escort-geometry} 与小节 \ref{subsec:conclusion-atomic-euler-ghost-primitive-core-zerotemp} 已把长度 \(2\) 原子 Euler 因子的 primitive/ghost/finite-part surgery 写成显式公式，而小节 \ref{subsec:conclusion-binfold-escort-blackwell-geometry-risk} 又已把 bin-fold escort 温度族压缩为一条两态 Bernoulli--logistic 骨架之后，还可把完成化截面、有限分辨率 Euler 包的低阶 jet、零温地基熵与 bulk 共振的剩余信息进一步闭合为如下结构结论。其核心事实是：标准完成化切片恰好把全部满足 \(\ell=2E\) 的 balanced atom 压扁为规范化核，而低于 \(F_{m+2}\) 阶的 canonical jet 又与 wrapped Poisson 热核完全重合；于是，``低阶 jet \(+\) normalized completion'' 构成了一类真正的双盲审计接口。与此同时，同一 balanced atom 在 primitive 与 Abel 有限部上仍留下可精确审计的残余，而零温地基层的微正则熵正好等于四态地基 SFT 拓扑熵的一半；再结合严格正的稳态电流，便可看出完成化的 involutive 自对偶与地基过程的 detailed balance 根本不是同一层级的刚性。

\begin{theorem}[完成化截面对平衡原子的常数化失明]\label{thm:conclusion-completed-slice-balanced-atoms-become-constants}
\leanverified{paper\_conclusion\_completed\_slice\_balanced\_atoms\_become\_constants}
设二变量动力学行列式 admit 有限原子分解
\[
\Delta(z,u)
=
\Bigl(\prod_{j=1}^{J}(1-u^{E_j}z^{\ell_j})^{m_j}\Bigr)\Delta_{\mathrm{core}}(z,u),
\qquad
E_j,\ell_j,m_j\in\ZZ_{>0}.
\]
对任意 \(L>1\)，定义完成化截面
\[
z_L(s):=L^{-s},
\qquad
u_L(s):=L^{2s-1},
\qquad
\Xi_{\Delta,L}(s):=\Delta(z_L(s),u_L(s)).
\]
则有严格恒等式
\[
\Xi_{\Delta,L}(s)
=
\Bigl(\prod_{j=1}^{J}\bigl(1-L^{-E_j}L^{(2E_j-\ell_j)s}\bigr)^{m_j}\Bigr)
\Xi_{\Delta_{\mathrm{core}},L}(s).
\]
特别地，凡满足
\[
\ell_j=2E_j
\]
的原子，在完成化截面上都退化为与 \(s\) 无关的常数因子
\[
(1-L^{-E_j})^{m_j}.
\]
\end{theorem}
\begin{proof}
对每个原子因子，直接代入完成化参数即可得到
\[
u_L(s)^{E_j}z_L(s)^{\ell_j}
=
L^{E_j(2s-1)}L^{-\ell_j s}
=
L^{-E_j}L^{(2E_j-\ell_j)s}.
\]
故
\[
(1-u_L(s)^{E_j}z_L(s)^{\ell_j})^{m_j}
=
\bigl(1-L^{-E_j}L^{(2E_j-\ell_j)s}\bigr)^{m_j}.
\]
逐项相乘即得所述乘积公式。若 \(\ell_j=2E_j\)，则 \(s\)-依赖完全消失，只剩常数因子 \((1-L^{-E_j})^{m_j}\)。证毕。
\end{proof}

\begin{theorem}[normalized completion 的 balanced atom 核刻画]\label{thm:conclusion-completed-slice-normalized-kernel-balanced-atoms}
\leanverified{paper\_conclusion\_completed\_slice\_normalized\_kernel\_balanced\_atoms}
沿用定理 \ref{thm:conclusion-completed-slice-balanced-atoms-become-constants} 的记号。对任意 \(L>1\)，定义规范化完成化截面
\[
\mathfrak C_L[\Delta](s)
:=
\frac{\Xi_{\Delta,L}(s)}{\Xi_{\Delta,L}(1/2)}.
\]
则以下两条严格等价：
\begin{enumerate}
  \item 对所有 \(L>1\) 与所有 \(s\)，都有
  \[
  \mathfrak C_L[\Delta](s)=\mathfrak C_L[\Delta_{\mathrm{core}}](s);
  \]
  \item 对每一个原子因子都成立
  \[
  \ell_j=2E_j.
  \]
\end{enumerate}
换言之，balanced atom 恰构成 normalized completion 的全部原子核；除它们以外，不存在别的有限原子在标准完成化归一化后能够从一变量截面中完全隐去。
\end{theorem}
\begin{proof}
由定理 \ref{thm:conclusion-completed-slice-balanced-atoms-become-constants}，
\[
\frac{\mathfrak C_L[\Delta](s)}{\mathfrak C_L[\Delta_{\mathrm{core}}](s)}
=
\prod_{j=1}^{J}
\left(
\frac{1-L^{-E_j}L^{(2E_j-\ell_j)s}}
{1-L^{-E_j}L^{(2E_j-\ell_j)/2}}
\right)^{m_j}.
\]
若对每个 \(j\) 都有 \(\ell_j=2E_j\)，则上式每一项都等于 \(1\)，故 \(\mathfrak C_L[\Delta](s)=\mathfrak C_L[\Delta_{\mathrm{core}}](s)\)。

反之，设存在某个指标 \(j\) 使
\[
\beta_j:=2E_j-\ell_j\neq 0.
\]
若某个 \(\beta_j>0\)，取
\[
\beta_{\max}:=\max_j \beta_j>0.
\]
则当实变量 \(s\to+\infty\) 时，所有满足 \(\beta_j=\beta_{\max}\) 的因子模长都满足
\[
\left|
\frac{1-L^{-E_j}L^{\beta_j s}}
{1-L^{-E_j}L^{\beta_j/2}}
\right|^{m_j}
\asymp
c_j(L)\,L^{m_j\beta_{\max}s},
\]
因而整个乘积的模长按
\[
L^{\beta_{\max}s\sum_{\beta_j=\beta_{\max}}m_j}
\]
指数增长，不可能恒等于 \(1\)。若不存在正的 \(\beta_j\)，则必有某个 \(\beta_j<0\)；对此取
\[
\beta_{\min}:=\min_j\beta_j<0,
\]
令 \(s\to-\infty\) 即得同样的指数增长矛盾。故不可能存在 \(\beta_j\neq 0\)，于是对所有 \(j\) 都必须有 \(\ell_j=2E_j\)。证毕。
\end{proof}

\begin{theorem}[真实输入 \texorpdfstring{$40$}{40} 状态核的完成化原子失明与 finite-part 残余]\label{thm:conclusion-realinput40-completed-slice-atomblind-finitepart-visible}
沿用附录 \ref{app:real-input-40-zeta-u} 的记号。对真实输入 \(40\) 状态核，
\[
\Delta(z,u)=(1-uz^2)\,F(z,u).
\]
则对任意 \(L>1\)，标准完成化截面满足
\[
\Xi_{\Delta,L}(s)=(1-L^{-1})\,\Xi_{F,L}(s).
\]
因而在一切分母非零的点上有
\[
\frac{\Xi_{\Delta,L}(s)}{\Xi_{\Delta,L}(1/2)}
=
\frac{\Xi_{F,L}(s)}{\Xi_{F,L}(1/2)}.
\]
换言之，完成化一变量对象的零点、重数、对数导数与一切标准化后的 \(s\)-依赖，均完全看不见显式 primitive 原子 \((1-uz^2)\)。

然而同一原子在二变量 primitive/finite-part 接口上仍留下严格残余：
\[
P(z,u)-P_{\mathrm{core}}(z,u)=uz^2,
\qquad
\log\mathfrak M(u)-\log\mathfrak M_{\mathrm{core}}(u)=z_\star(u)^2.
\]
因此，标准完成化截面与 Abel 有限部之间存在严格的信息分裂：前者对平衡原子全失明，而后者精确保留其 residual displacement。
\end{theorem}
\begin{proof}
把定理 \ref{thm:conclusion-completed-slice-balanced-atoms-become-constants} 应用于单原子
\[
(E,\ell,m)=(1,2,1)
\]
即可得到
\[
1-u_L(s)z_L(s)^2
=
1-L^{2s-1}L^{-2s}
=
1-L^{-1},
\]
从而
\[
\Xi_{\Delta,L}(s)=(1-L^{-1})\,\Xi_{F,L}(s).
\]
第二条比例恒等式只是把同一常数因子在标准化中约去。

另一方面，命题 \ref{prop:atomic-2-orbit-split-logM} 已给出
\[
P(z,u)=uz^2+P_{\mathrm{core}}(z,u),
\qquad
\log\mathfrak M(u)=\log\mathfrak M_{\mathrm{core}}(u)+z_\star(u)^2.
\]
故同一原子在 primitive 与 finite-part 层保留完全显式的差额。证毕。
\end{proof}

\begin{theorem}[有限分辨率 Euler 包的低阶 jet 与 normalized completion 双盲]\label{thm:conclusion-zeckendorf-euler-lowjet-completion-doubleblind}
固定分辨率 \(m\)，并记
\[
F:=F_{m+2}.
\]
设 \(\mu_{m,t}\) 为 canonical Euler 包的 factorial section。则以下两类有限审计同时失明：
\begin{enumerate}
  \item 任何只依赖于 Fourier 侧低阶 jet 数据
  \[
  \left\{
  \frac{\dd^n}{\dd t^n}\widehat\mu_{m,t}(k)\Big|_{t=0}
  :
  0\le n\le F-1,\ k\in G_m
  \right\};
  \]
  \item 任何只依赖于 normalized completion 数据
  \[
  \mathfrak C_L[\Delta](s)=\frac{\Xi_{\Delta,L}(s)}{\Xi_{\Delta,L}(1/2)}.
  \]
\end{enumerate}
更精确地，\(\mu_{m,t}\) 与 wrapped Poisson 热核在 \(t=0\) 处具有完全相同的 \((F-1)\)-阶 Taylor germ，而一切 balanced atom 在 normalized completion 中都完全不可见。因而任何仅允许 ``低于 \(F\) 阶 jet \(+\) normalized completion'' 这两类输入的有限审计器，都无法同时检出有限分辨率非半群缺陷与 balanced atom 的存在；第一可见见证必须至少越过这两道门槛之一。
\end{theorem}
\begin{proof}
命题 \ref{prop:conclusion-zeckendorf-euler-fminusone-jet-rigidity} 已证明：对每个 \(k\in G_m\)，都有
\[
\widehat\mu_{m,t}(k)
=
e^{\,t(e^{-2\pi ik/F}-1)}+O(t^F),
\]
故
\[
\frac{\dd^n}{\dd t^n}\widehat\mu_{m,t}(k)\Big|_{t=0}
=
\bigl(e^{-2\pi ik/F}-1\bigr)^n
\qquad
(0\le n\le F-1),
\]
与 wrapped Poisson 热核完全一致。另一方面，定理 \ref{thm:conclusion-completed-slice-normalized-kernel-balanced-atoms} 已表明：全部 balanced atom 恰落在 normalized completion 的原子核中。于是，凡仅因子化经由上述两类数据的有限审计器，都既不能在低于 \(F\) 阶的 jet 层看见非半群缺陷，也不能在 completion 层看见 balanced atom。证毕。
\end{proof}

\leanverified{paper\_conclusion\_realinput40\_ground\_entropy\_half\_topological\_entropy}
\begin{theorem}[极端饱和微正则熵与零温四态地基 SFT 的半熵恒等式]\label{thm:conclusion-realinput40-ground-entropy-half-topological-entropy}
设 \(A_{n,k_{\max}(n)}\) 为输出位势下达到最大能量密度的微观计数，并定义
\[
h_{\mathrm{ground}}
:=
\lim_{n\to\infty}\frac1n\log A_{n,k_{\max}(n)}.
\]
又令 \(B_\infty\) 为零温缩放极限中的四状态非负整数矩阵，\(\rho(B_\infty)\) 为其 Perron 根。则
\[
h_{\mathrm{ground}}
=
\frac12\log \rho(B_\infty)
=
\frac12\,h_{\mathrm{top}}(\Sigma_{B_\infty}).
\]
换言之，极端饱和微正则熵恰等于零温地基 SFT 拓扑熵的一半。
\end{theorem}
\begin{proof}
由推论 \ref{cor:real-input-40-ground-entropy}，
\[
h_{\mathrm{ground}}=\log c.
\]
另一方面，定理 \ref{thm:xi-time-part73-perron-vs-atomic-sqrtu-separation} 给出
\[
\rho(B_\infty)=c^2.
\]
于是
\[
h_{\mathrm{ground}}=\log c=\frac12\log c^2=\frac12\log\rho(B_\infty).
\]
对有限型移位 \(\Sigma_{B_\infty}\)，其拓扑熵等于 \(\log\rho(B_\infty)\)。代入即得
\[
h_{\mathrm{ground}}=\frac12\,h_{\mathrm{top}}(\Sigma_{B_\infty}).
\]
证毕。
\end{proof}

另一方面，bin-fold escort 温度族的末位读出塌缩与 logistic 信息几何已经在定理 \ref{thm:conclusion-binfold-escort-blackwell-equivalence}、定理 \ref{thm:killo-fold-bin-escort-renyi-logistic-geometry} 与定理 \ref{thm:conclusion-escort-fisher-rao-geodesic-reparametrization} 中被完全显式化：整个实验族在指数精度上等价于单个 Bernoulli\((\theta(q))\) 实验，而其 KL、Hellinger、Chernoff 与 Fisher 结构全部退化为同一条 logistic 曲线。把这条一比特骨架与 bulk 共振常数 \(C_\varphi\) 的碰撞缺口下界并置，即可得到如下不可消除的二阶均匀性障碍。

\begin{theorem}[bulk 共振导致的 Rényi--\texorpdfstring{$2$}{2} 均匀性常数缺口]\label{thm:conclusion-binfold-renyi2-uniformity-obstruction-from-resonance}
令 \(p_m\) 为折叠碰撞在 \(\ZZ/F_{m+2}\ZZ\) 上的归一化分布，\(u_m\) 为同一群上的均匀分布，并记
\[
C_\varphi:=\prod_{n=2}^{\infty}\bigl|\cos(\pi\varphi^{-n})\bigr|.
\]
则存在严格的二阶均匀性障碍
\[
\liminf_{m\to\infty} D_2(p_m\Vert u_m)\ge \log(1+2C_\varphi^2)>0.
\]
并且
\[
\liminf_{m\to\infty}
F_{m+2}\,\|p_m-u_m\|_2^2
\ge
2C_\varphi^2,
\]
从而
\[
\|p_m-u_m\|_2
\ge
\frac{\sqrt2\,C_\varphi+o(1)}{\sqrt{F_{m+2}}},
\qquad
\|p_m-u_m\|_{\mathrm{TV}}
\ge
\frac{C_\varphi+o(1)}{\sqrt{2F_{m+2}}}.
\]
因此，bulk 共振造成的偏离不是可由高频误差吸收的尾部瑕疵，而是在 Rényi--\(2\)、\(\ell^2\) 与全变差三种量纲上同时保留的常数量级缺口。
\end{theorem}
\begin{proof}
定理 \ref{thm:killo-fold-renyi2-uniformity-gap} 已给出
\[
\liminf_{m\to\infty} D_2(p_m\Vert u_m)\ge \log(1+2C_\varphi^2),
\]
以及
\[
\liminf_{m\to\infty}F_{m+2}\!\left(\mathsf{Col}_m-\frac1{F_{m+2}}\right)\ge 2C_\varphi^2.
\]
再由命题 \ref{prop:fold-bin-chi2-col}，
\[
\|p_m-u_m\|_2^2=\mathsf{Col}_m-\frac1{F_{m+2}}.
\]
于是
\[
\liminf_{m\to\infty}F_{m+2}\,\|p_m-u_m\|_2^2\ge 2C_\varphi^2.
\]
取平方根便得到
\[
\|p_m-u_m\|_2\ge \frac{\sqrt2\,C_\varphi+o(1)}{\sqrt{F_{m+2}}}.
\]
最后利用
\[
\|x\|_1\ge \|x\|_2,
\qquad
\|p_m-u_m\|_{\mathrm{TV}}=\frac12\|p_m-u_m\|_1,
\]
即可推出
\[
\|p_m-u_m\|_{\mathrm{TV}}
\ge
\frac{C_\varphi+o(1)}{\sqrt{2F_{m+2}}}.
\]
证毕。
\end{proof}

\begin{theorem}[零温四态地基 SFT 的 Parry 链闭式]\label{thm:conclusion-realinput40-zero-temp-parry-chain-explicit}
\leanverified{paper\_conclusion\_realinput40\_zero\_temp\_parry\_chain\_explicit}
设
\[
B_\infty=
\begin{pmatrix}
0&1&1&0\\
0&0&1&0\\
0&0&3&1\\
1&0&0&3
\end{pmatrix},
\qquad
\rho:=\rho(B_\infty),
\]
其中 \(\rho\) 满足
\[
\rho^4-6\rho^3+9\rho^2-\rho-1=0.
\]
则其 Parry 极大熵测度对应的转移矩阵可写成
\[
P_\infty=
\begin{pmatrix}
0 & \dfrac{1}{\rho+1} & \dfrac{\rho}{\rho+1} & 0\\[4pt]
0 & 0 & 1 & 0\\[4pt]
0 & 0 & \dfrac{3}{\rho} & \dfrac{\rho-3}{\rho}\\[4pt]
\dfrac{\rho-3}{\rho} & 0 & 0 & \dfrac{3}{\rho}
\end{pmatrix},
\]
其平稳分布为
\[
\pi_\infty
=
\frac{1}{3\rho^2+\rho-6}
\bigl(
\rho^2-2\rho-3,\;
\rho-3,\;
\rho^2+\rho,\;
\rho^2+\rho
\bigr).
\]
特别地，
\[
\pi_{\infty,3}=\pi_{\infty,4},
\qquad
P_{\infty,23}=1.
\]
\end{theorem}
\begin{proof}
定理 \ref{thm:killo-real-input-40-zero-temp-ground-sft-parry-closed-form} 已给出 Parry 转移矩阵的显式公式，即上式所示 \(P_\infty\)。

现求平稳分布。设 \(\pi=(\pi_1,\pi_2,\pi_3,\pi_4)\) 满足 \(\pi P_\infty=\pi\)。由第一、第二与第四列的平衡方程得到
\[
\pi_1=\frac{\rho-3}{\rho}\pi_4,
\qquad
\pi_2=\frac{1}{\rho+1}\pi_1,
\qquad
\pi_4=\frac{\rho-3}{\rho}\pi_3+\frac{3}{\rho}\pi_4.
\]
最后一式等价于
\[
(\rho-3)\pi_4=(\rho-3)\pi_3,
\]
而 \(\rho>3\)，故
\[
\pi_3=\pi_4.
\]
令 \(\pi_4=x\)。则
\[
\pi_3=x,
\qquad
\pi_1=\frac{\rho-3}{\rho}x,
\qquad
\pi_2=\frac{\rho-3}{\rho(\rho+1)}x.
\]
由归一化条件 \(\pi_1+\pi_2+\pi_3+\pi_4=1\)，得到
\[
1=x\left(
2+\frac{\rho-3}{\rho}+\frac{\rho-3}{\rho(\rho+1)}
\right)
=
x\,\frac{3\rho^2+\rho-6}{\rho(\rho+1)}.
\]
故
\[
x=\frac{\rho(\rho+1)}{3\rho^2+\rho-6}.
\]
代回并化简，便得到
\[
\pi_1=\frac{(\rho+1)(\rho-3)}{3\rho^2+\rho-6}
=
\frac{\rho^2-2\rho-3}{3\rho^2+\rho-6},
\]
\[
\pi_2=\frac{\rho-3}{3\rho^2+\rho-6},
\qquad
\pi_3=\pi_4=\frac{\rho^2+\rho}{3\rho^2+\rho-6}.
\]
证毕。
\end{proof}

\begin{theorem}[零温四态地基支撑的不可逆性]\label{thm:conclusion-realinput40-ground-support-irreversibility}
\leanverified{paper\_conclusion\_realinput40\_ground\_support\_irreversibility}
不存在任何与定理 \ref{thm:conclusion-realinput40-zero-temp-parry-chain-explicit} 中 \(P_\infty\) 具有相同有向支撑的可逆 Markov 链。特别地，零温地基层的半熵凝聚并不是某个保持同支撑的 detailed-balance 投影。
\end{theorem}
\begin{proof}
设 \(Q\) 是一个与 \(P_\infty\) 具有相同有向支撑的 Markov 链，并假设它可逆。由于 \(Q\) 与 \(P_\infty\) 拥有相同的强连通有向图，\(Q\) 是有限不可约链，故其平稳分布 \(\pi^Q\) 处处严格正。

另一方面，由定理 \ref{thm:conclusion-realinput40-zero-temp-parry-chain-explicit}，
\[
P_{\infty,23}=1,
\qquad
P_{\infty,32}=0.
\]
同支撑意味着
\[
Q_{23}>0,
\qquad
Q_{32}=0.
\]
若 \(Q\) 可逆，则对其平稳分布必须满足 detailed balance：
\[
\pi^Q_2Q_{23}=\pi^Q_3Q_{32}.
\]
右端等于 \(0\)，而左端因 \(\pi^Q_2>0\) 且 \(Q_{23}>0\) 而严格为正，矛盾。故不存在这样的可逆链。证毕。
\end{proof}

\begin{corollary}[完成化自对偶不蕴含零温地基的 detailed balance]\label{cor:conclusion-realinput40-selfdual-completion-not-detailed-balance}
\leanverified{paper\_conclusion\_realinput40\_selfdual\_completion\_not\_detailed\_balance}
标准完成化截面满足严格函数方程
\[
\Xi_{\Delta,L}(s)=\Xi_{\Delta,L}(1-s),
\]
但零温地基层的 Parry 链并不因此可逆。事实上，对定理 \ref{thm:conclusion-realinput40-zero-temp-parry-chain-explicit} 中的平稳分布与转移矩阵，有
\[
J_{23}
:=
\pi_{\infty,2}P_{\infty,23}-\pi_{\infty,3}P_{\infty,32}
=
\frac{\rho-3}{3\rho^2+\rho-6}
>0.
\]
因此，完成化自对偶是截面上的 involutive 谱对称，而不是零温地基动力学上的 detailed balance。
\end{corollary}
\begin{proof}
完成化函数方程来自标准完成化截面的定义。另一方面，由定理 \ref{thm:conclusion-realinput40-zero-temp-parry-chain-explicit}，
\[
P_{\infty,23}=1,
\qquad
P_{\infty,32}=0,
\qquad
\pi_{\infty,2}=\frac{\rho-3}{3\rho^2+\rho-6}>0.
\]
故
\[
J_{23}
=
\pi_{\infty,2}P_{\infty,23}-\pi_{\infty,3}P_{\infty,32}
=
\pi_{\infty,2}
=
\frac{\rho-3}{3\rho^2+\rho-6}
>0.
\]
因此 \(\pi_{\infty,i}P_{\infty,ij}=\pi_{\infty,j}P_{\infty,ji}\) 不可能对全部 \(i,j\) 成立，Parry 链不是可逆 Markov 链。证毕。
\end{proof}

综上，标准完成化切片不只是把 balanced atom 压扁为与 \(s\) 无关的常数，而且在归一化后恰好以这些平衡原子作为全部 completion kernel；再与 \((F_{m+2}-1)\) 阶 Poisson jet 刚性并置，便得到一条真正的双盲审计链：低阶 jet 看不见有限分辨率非半群缺陷，normalized completion 看不见 balanced atom，而 primitive 与 Abel 有限部却继续精确保留同一原子的横向位移。与此同时，零温饱和微正则熵恰等于四态地基 SFT 拓扑熵的一半，bulk 共振又在 Rényi--\(2\) 与 \(\ell^2\)/总变差层面保留严格正的均匀性缺口；最后，完成化函数方程与同支撑不可逆的零温 Parry 流并存，说明 involutive 自对偶与 detailed balance 属于两种彼此独立的刚性。

\endinput
